\section{Limitations}
\label{sec:limitations}

CRA5 reportedly trained on 1979--2021 and tested on 2022--2023~\cite{han2024};
we did not audit its checkpoint log. Our 2018 test overlaps CRA5 training,
while one frozen 2022 test, one CONUS region, one neural codec and few
correctors do not establish broad generalisation.

Matched mean rates compare different global/regional coding tasks. Frame-rate
coefficients of variation are 0.95\% for CRA5 and 3.70\% for MPCA-only.
Payloads exclude shared bases and correction weights; inference adds cost.
Derived \texttt{tcwv} includes derivation error. No pointwise or per-block
improvement is guaranteed.

Seven-day bootstrap intervals condition on checkpoint, year and grid; they do
not remove all temporal or spatial dependence. Roughness covaries with elevation
and geography; western lowland bins have only 29--31 cells, while mountain
pressure levels may be below ground. Terrain associations are not causal.
One dilation pattern/seed cannot exclude broader context or optimisation effects;
unequal group sizes and error magnitudes affect the shared residual-MSE loss.
